\documentclass[mathserif, 10pt]{beamer}
\usepackage[utf8]{inputenc}
\usepackage{amsmath, amsfonts}
\usepackage{appendixnumberbeamer}


\title{Tutorial: python y flavio}
\subtitle{Jorge Alda,\\ Universidad de Zaragoza/CAPA\hspace{4em} \texttt{jalda@unizar.es} }
\author[Jorge Alda]{Slides Example}

\date{Mayo 2022}



\usetheme{Zaragoza}
\usecolortheme{Unizar}
\titlepagelogoA{\includegraphics[width=6cm]{logos/CAPA.png}}
\titlepagelogoB{\includegraphics[width=4cm]{logos/dftuz2.png}}


\newcommand\colorcite[1]{{\scriptsize\color{blue}#1}}

\begin{document}
\begin{frame}[noframenumbering,plain]

    \titlepage

\end{frame}


\begin{frame}\frametitle{venv}
    Un ambiente virtual (venv) permite ``encapsular'' los paquetes de \texttt{python} que necesitas de forma aislada del resto del sistema operativo.

    ~

    Para crear un venv en el directorio actual:
    \begin{block}{Linux}
        \texttt{python3 -m venv envname}
    \end{block}
    \begin{block}{Windows}
        \texttt{C:\textbackslash Path\textbackslash to\textbackslash python -m venv envname}
    \end{block}

    Esto crea un subdirectorio \texttt{envname} con los archivos necesarios para usar el \texttt{venv}.
\end{frame}

\begin{frame}\frametitle{venv}

    Ahora hay que activar el texttt{venv}
    \begin{block}{Linux}
        \texttt{source envname/bin/activate}
    \end{block}
    \begin{block}{Windows}
        \texttt{envname\textbackslash Scripts\textbackslash Activate.ps1}
    \end{block}

    Comprueba que estás dentro del \texttt{venv} usando el comando \texttt{python -m pip list}.

    Para salir del \texttt{venv}, usa el comando
    \begin{block}{Windows \& Linux}
        deactivate
    \end{block}
\end{frame}

\begin{frame}\frametitle{Preparando el venv}
    Activa el \texttt{venv}, e instala los siguientes paquetes usando \texttt{pip}:
    \begin{itemize}
        \item numpy
        \item scipy
        \item matplotlib
\item jupyter
        \item flavio
    \end{itemize}
\end{frame}

\begin{frame}\frametitle{git y GitHub}

    \begin{itemize}
        \item \texttt{git} es un sistema de control de versiones. Permite mantener una historia de todos los cambios de un proyecto de código y trabajar en paralelo con varias versiones (ramas).
              \begin{itemize}
                  \item Cliente de línea de comandos en \url{https://git-scm.com/downloads}
                  \item También hay clientes gráficos o integrados en IDE (VSCode).
              \end{itemize}
        \item \texttt{GitHub} es un servidor de git. Permite mantener una copia ``en la nube'' del proyecto y su historia, editar de forma colaborativa, y compartir el código (open software).
              \begin{itemize}
                  \item Crear una cuenta en \url{https://github.com/}
                  \item Cliente de línea de comandos en \url{https://github.com/cli/cli}
              \end{itemize}
        \item Conecta \texttt{git} con la cuenta de \texttt{GitHub} con el comando \texttt{gh auth login} y sigue las instrucciones (cuenta de \texttt{GitHub} básica, protocolo HTTPS, autenticación usando browser).
    \end{itemize}
\end{frame}

\begin{frame}\frametitle{Empezando con git}
    \begin{enumerate}
        \item Entra en github.com con tu cuenta, pulsa el botón + en la esquina superior derecha y crea un nuevo repositorio. El nombre del repositorio es test-github, y añade un README.
        \item En la página del repositorio, pulsa el botón verde ``Code'' y copia la dirección. En la consola de comandos, navega al directorio donde quieras crear tu proyecto y ejecuta \texttt{git clone https://github.com/your-name/test-git.git}. El archivo README.md se habrá copiado a tu ordenador.
\end{enumerate}
\end{frame}

\begin{frame}\frametitle{Haciendo cambios en git}
\begin{enumerate}\setcounter{enumi}{2}
\item Crea un archivo test.txt en el directorio y guárdalo.
\item Ejecuta \texttt{git status}. Verás que el archivo test.txt ha sido modificado desde la última entrada en la historia del repositorio.
\item Ejecuta \texttt{git add test.txt}. Esto añade el archivo a la lista de preparación (stage), pero aún no a la historia.
\item Vuelve a ejecutar \texttt{git status}, el archivo está en el stage.
\item Una vez que hayas preparado todos los archivos modificados, puedes añadir una nueva entrada (commit) a la historia local con el comando \texttt{git commit -m "Descripción del commit"} (intenta ser descriptivo).
\item Puedes ver todos los commits locales con \texttt{git log}.
\end{enumerate}
    

\end{frame}

\begin{frame}\frametitle{Sincronizando con GitHub}
    \begin{enumerate}\setcounter{enumi}{8}
        \item Para subir los nuevos commits a GitHub, usa \texttt{git push}.
        \item En GitHub puedes ver la historia del repositorio.
        \item Si hay cambios en GitHub que no están en el ordenador, puedes incorporarlos con \texttt{git pull}. \textbf{Importante!:} Haz pull antes de commit+push.
    \end{enumerate}

    GitHub tiene otras funcionalidades avanzadas (issues, pull requests,...) que no utilizaremos de momento.

\end{frame}

\begin{frame}\frametitle{Ramas}
    Hasta ahora, solo hemos usado historias lineales: cada commit solo tiene un predecesor o un sucesor. Una historia con varias ramas sirve, por ejemplo, para probar nuevos conceptos mientras se mantiene (y desarrolla) la versión original.

    \begin{enumerate}
        \item Crea una rama con el comando \texttt{git branch rama1}. La rama rama1 tiene la misma historia hasta este punto que la rama main.
        \item Puedes ver todas las ramas del repositorio (locales+remotas) con \texttt{git branch -a}
        \item Para cambiar la rama activa, usa \texttt{git checkout rama1}.
    \end{enumerate}

\end{frame}

\end{document}